
\documentclass[11pt]{article}

\usepackage{amsmath,amssymb,amsthm}
\usepackage{geometry}
\usepackage{hyperref}
\usepackage{physics}
\usepackage{mathrsfs}
\usepackage{enumitem}

\geometry{margin=1in}

\title{\textbf{The Ehlers--Pirani--Schild Construction and the Geometric Obstruction to Classical Magnetism}}

\author{}

\date{}

\begin{document}

\maketitle

\begin{abstract}
The Ehlers--Pirani--Schild (EPS) framework reconstructs spacetime geometry from physically observable structures: the trajectories of free massive particles and light rays. In this paper we show that, within this framework, classical electromagnetism cannot generate intrinsic magnetic properties of matter. We argue that the absence of magnetism in classical statistical mechanics is not accidental but a consequence of the geometric constraints imposed by projective and conformal structures. The analysis reveals a deep connection between the no–interaction theorem, the flatness of classical phase space, and the impossibility of classical magnetism. Magnetism emerges only when the geometric foundations of classical spacetime are transcended, as in quantum theory.
\end{abstract}

\tableofcontents

%===================================================
\section{Introduction}

Classical electrodynamics is often regarded as a complete theory of electromagnetic phenomena in the macroscopic world. However, when coupled to classical statistical mechanics, it famously fails to account for magnetism in equilibrium. This failure is not merely technical; it reflects a deep incompatibility between magnetism and the geometric foundations of classical physics.

The Ehlers–Pirani–Schild (EPS) program provides a natural setting to understand this incompatibility. By reconstructing spacetime geometry from primitive physical structures—light rays and free-fall trajectories—it exposes which interactions are compatible with relativistic covariance and which are not.

In this paper, we argue that magnetism is excluded at a foundational level within the EPS framework. The obstruction is geometric rather than dynamical and is closely related to the no–interaction theorem for relativistic particle systems.

---

\section{The EPS Program: Geometry from Physics}

\subsection{Primitive Structures}

The EPS approach begins with two empirically defined families of curves:

\begin{itemize}
\item \textbf{Light rays}: null trajectories defined operationally by the propagation of light.
\item \textbf{Free-fall worldlines}: trajectories of freely falling massive test particles.
\end{itemize}

No metric is assumed a priori.

---

\subsection{Conformal and Projective Structures}

From light rays one reconstructs a \emph{conformal structure} $[g]$, defining null cones but not lengths.

From free-fall trajectories one reconstructs a \emph{projective structure}, i.e., an equivalence class of affine connections sharing the same unparameterized geodesics.

A spacetime admits a pseudo-Riemannian metric $g$ if and only if the conformal and projective structures are compatible.

---

\section{Electromagnetism in the EPS Framework}

\subsection{Charged Particles and Non-Geodesic Motion}

Charged particles do not follow geodesics of the Levi-Civita connection:
\[
\frac{D u^\mu}{d\tau} = \frac{q}{m} F^\mu{}_{\nu} u^\nu.
\]

Thus electromagnetic forces do not define a new projective structure; they merely deform trajectories relative to the gravitational background.

---

\subsection{Gauge Fields as External Structures}

The electromagnetic potential defines a connection on a $U(1)$ bundle over spacetime, but it does not modify the affine structure of spacetime itself.

Hence:
\begin{itemize}
\item Gravity shapes spacetime geometry.
\item Electromagnetism acts within a fixed geometry.
\end{itemize}

This asymmetry is central to the absence of magnetism in classical equilibrium.

---

\section{Phase Space Geometry and Canonical Structure}

\subsection{Cotangent Bundle Structure}

Classical phase space is the cotangent bundle $T^*M$ equipped with its canonical symplectic form
\[
\omega = dp_\mu \wedge dx^\mu.
\]

Minimal electromagnetic coupling modifies the Hamiltonian but leaves $\omega$ unchanged.

---

\subsection{Canonical Flatness and Gauge Redundancy}

Because the symplectic structure is exact, the transformation
\[
p_\mu \mapsto p_\mu - q A_\mu
\]
is canonical.

Thus the electromagnetic field introduces no intrinsic curvature into phase space.

This fact underlies the Bohr–van Leeuwen theorem.

---

\section{No–Interaction Theorem and EPS Geometry}

The Currie–Jordan–Sudarshan theorem states that relativistic particle systems with canonical coordinates and Poincaré invariance admit no nontrivial interactions.

In EPS language:
\begin{itemize}
\item Projective structure fixes free motion.
\item Any interaction that preserves relativistic covariance must be geometrizable.
\item Electromagnetism fails to modify the projective structure.
\end{itemize}

Therefore it cannot induce collective effects such as magnetization.

---

\section{Why Magnetism Is Excluded}

Magnetism requires:
\begin{enumerate}
\item persistent microscopic currents,
\item energetic preference for alignment,
\item broken time-reversal symmetry at equilibrium.
\end{enumerate}

EPS geometry forbids all three:
\begin{itemize}
\item canonical measures are time-reversal symmetric,
\item no intrinsic spin exists,
\item no curvature is induced in phase space.
\end{itemize}

Hence classical magnetism is geometrically impossible.

---

\section{Quantum Contrast}

Quantum theory alters the geometric foundations:

\begin{itemize}
\item Phase space becomes noncommutative.
\item Gauge fields acquire holonomy.
\item Berry curvature replaces symplectic flatness.
\item Spin emerges as an intrinsic representation of the Lorentz group.
\end{itemize}

These features evade the EPS constraints.

---

\section{Conceptual Synthesis}

The EPS program shows that classical spacetime geometry severely restricts admissible dynamics. Magnetism, being intrinsically geometric and topological, lies beyond those restrictions.

Thus the absence of magnetism is not a deficiency of classical electrodynamics but a theorem about classical spacetime itself.

---

\section{Conclusion}

We have shown that within the Ehlers–Pirani–Schild framework, magnetism is excluded on purely geometric grounds. Classical electrodynamics operates within a phase space whose structure forbids intrinsic magnetic response. Only by abandoning classical spacetime geometry—through quantum mechanics—does magnetism become possible.

---

\section*{References}

\begin{enumerate}
\item J. Ehlers, F. A. E. Pirani, A. Schild, \emph{The Geometry of Free Fall and Light Propagation}, in \emph{General Relativity}, Clarendon Press (1972).
\item J. Ehlers, \emph{Survey of General Relativity Theory}, in \emph{Relativity, Astrophysics and Cosmology}, Reidel (1973).
\item L. D. Landau, E. M. Lifshitz, \emph{The Classical Theory of Fields}.
\item R. Haag, \emph{Local Quantum Physics}.
\item J. Baez, J. P. Muniain, \emph{Gauge Fields, Knots and Gravity}.
\item H. Goldstein, C. Poole, J. Safko, \emph{Classical Mechanics}.
\end{enumerate}

\end{document}
